\cotefile{157-3}
\scanpage{342}
\mspage{41}

de l’âme est sans commencement et sans fin ! Cela signifie, par contre, libérer peu ou prou ce désir des vibrations parasites issues de l’égo qui le brouillent. Parmi celles-ci, il y a surtout celles qui émanent de la redoutable figure de \textit{Censeur}, gardien des répressions qui s’évertuent à cacher ce désir et à le désavouer. Avec chaque libération partielle, instaurant notre autonomie par rapport à l’une ou l’autre des figures archétypes, le Censeur s’amenuise, il perd de son autorité et de sa prestance. Et quand nous sommes libres enfin et de nos Parents, et de nos Conjoints, le Censeur n’est plus que l’ombre de lui-même ! Une fois la guerre des générations et la guerre du sexe terminée, la mort du Censeur ne saurait plus tarder...

Parmi les quatre pulsions amoureuses archétypes, nous attirant en le corps de la Mère, en le corps de l’Epouse, en le corps du Père et en celui de l’Epoux, deux sont incestueuses, et deux sont homosexuelles - et pourtant elles sont ! Et elles sont éternellement créatrices et éternellement belles. Une seule parmi les quatre (celle qui porte l’homme pubère vers l’Epouse, la femme nubile vers l’Epoux) est censée être recevable, ou ne serait-ce qu’\textit{exister} selon le Plan de Dieu, ou selon les lois de Nature. Mais il n’en est pas une seule qui ne soit en Dieu de toute éternité, et qu’Il n’ait insufflée dans l’âme dès sa création, pour animer son élan vers ses destinées éternelles.

\textit{La grande tâche de l’âme au temps de son errance, celle qui précède et prépare toutes les autres, c’est de se libérer de sa dépendance par rapport à ses Amants et par rapport à ses Amantes originels, et ce faisant, d’apprendre à aimer.} Et aimer dans la plénitude, c’est aussi prendre connaissance du désir éternel qui nous porte vers chacun d’eux, et assumer ce désir dans toute sa force et dans toute sa beauté. Quand cette tâche-là est remplie, l’âme est prête enfin pour entrer dans les splendeurs éternelles.

\noindent 17. Assumer le désir, ce n’est pas nécessairement l’assouvir

Assumer un désir charnel, ce n’est ni plus ni moins qu’en prendre connaissance au niveau conscient, et d’en sentir intensément la beauté. C’est aussi ``\textit{vivre}'' ce désir dans sa force et dans sa beauté originelle, le vivre comme il est vécu par Dieu en Lui-même, et comme il est vécu aussi par les couches profondes de l’âme, en certains moments privilégiés. Lorsque (comme il est presque toujours le cas) le désir se heurte à des interdits intériorisés par la psyché, \textit{assumer ce désir est un acte de vérité et de courage, un acte de fidélité}, dont la portée est d’autant plus grande que les interdits sont plus puissants. Dans la progression spirituelle de l’âme, ces actes-là ne jouent pas un rôle moins crucial que ceux qui tendent à la libérer de la vanité. Les uns tendent à instaurer une autonomie par rapport au ``\textit{Flatteur}'', les autres une autonomie non moins ardue par rapport au Censeur.

Il est important de réaliser qu’assumer, donc aussi ``vivre'' un désir, est un acte d’essence différente de l’acte d’assouvissement. Aucun des deux actes n’implique l’autre : nous pouvons vivre un désir, ``vibrer'' en lui dans une plénitude parfaite, sans jamais l’assouvir ``en actes'' ni même (si les circonstances sont adverses) souhaiter l’assouvir ainsi, charnellement. Inversement, même les désirs charnels ressentis comme parfaitement orthodoxes sont presque toujours assouvis sans être véritablement assumés. Il en est ainsi quand les comportements amoureux sont dominés par les poses conventionnelles (ou d’autres poses de nature plus personnelles), imposées par la sempiternelle immixition de l’égo dans les ébats amoureux.

% archive sheet number: 342
